\begin{recipe}{Blini's met vier toppings}
\kicker[Voorgerecht,Oost-Europees,Kerst]{Voorgerecht · Oost-Europees · kerst}
\index[register]{Blini's met vier toppings}
\index[register]{Blini's}
\index[register]{Gerookte zalm}
\index[register]{Rosbief}
\index[register]{Geitenkaas}
\index[register]{Boekweitmeel}
\index[register]{Gist}
\index[register]{Appel}
\index[register]{Dille}

\meta{45 minuten \dvd 8 stuks}

\lettrine{B}{lini's} zijn kleine, dikke pannenkoekjes uit Oost-Europa en een
handig hapje voor bij de borrel: je belegt ze koud, dus er komt geen kookwerk
meer bij zodra ze klaar zijn. Ze zijn ook een leuk voorgerecht met kerst.
Deze vier toppings, met zalm, met rosbief, een vegetarische met
geitenkaas en nog een zalmvariant met appel en dille, maken samen een
gevarieerd plankje.

\blockrule{Ingrediënten}
\begin{ingredients}
  \ing{100 g}{boekweitmeel}\index[register]{Boekweitmeel}
  \ing{150 ml}{lauwwarme melk (150\% hydratatie)}
  \ing{1 tl}{zout (ca.\ 5\%)}
  \ing{1}{ei}
  \ing{7 g}{droge gist}\index[register]{Gist}
  \ingb{zonnebloemolie, om te bakken}
\end{ingredients}

\medskip
\noindent\textit{Toppings (elk voor 2 blini's)}
\begin{ingredients}
  \ingb{\textbf{Zalm:} 2 plakjes gerookte zalm\index[register]{Gerookte zalm}, 2 tl
        kruidenroomkaas (bijv.\ Boursin)\index[register]{Roomkaas}, 3 takjes
        fijngesneden bieslook\index[register]{Bieslook}}
  \ingb{\textbf{Rosbief:} 2 plakjes rosbief (met peper of
        naturel)\index[register]{Rosbief}, 10 g rucola\index[register]{Rucola},
        2 tl groene pesto\index[register]{Pesto}, 2 tl geraspte Parmezaanse
        kaas\index[register]{Parmezaanse kaas}}
  \ingb{\textbf{Geitenkaas:} 2 schijfjes zachte
        geitenkaas\index[register]{Geitenkaas}, 2 zongedroogde tomaten (op
        olie)\index[register]{Zongedroogde tomaten}, 2 tl vloeibare honing,
        1 takje verse tijm}
  \ingb{\textbf{Zalm, appel \& dille:} 2 plakjes gerookte
        zalm\index[register]{Gerookte zalm}, 2 tl zure
        room\index[register]{Zure room}, 2 takjes fijngesneden
        bieslook\index[register]{Bieslook}, 1/8 appel in piepkleine
        blokjes\index[register]{Appel}, dille\index[register]{Dille}}
\end{ingredients}

\blockrule{Bereiding}
\tip{Maak de vier toppings vlak voor het serveren op: de blini's worden
anders snel zacht.}
\begin{steps}
  \item Los de gist op in de lauwwarme melk en laat 5 minuten staan tot het
        een beetje schuimt. Splits het ei en zet het eiwit apart in de
        koelkast.
  \item Meng het boekweitmeel en zout met de eidooier en roer er de gist en
        melk doorheen tot een glad beslag. Zet 20 minuten opzij om te rijzen.
  \item Klop het eiwit stijf en spatel het voorzichtig door het gerezen
        beslag. Doe het beslag in een knijpfles (zie basisapparatuur).
  \item Verhit een beetje zonnebloemolie op middelhoog vuur in een koekenpan
        en knijp er kleine cirkeltjes beslag direct in. Bak de blini's 1 tot
        2 minuten per kant tot ze goudbruin zijn.
  \item Leg de blini's op een bord of plank.
  \item Voor de zalmvariant: schep op twee blini's een beetje kruidenroomkaas
        en leg daar een plakje gerookte zalm overheen. Werk af met de
        fijngesneden bieslook.
  \item Voor de rosbiefvariant: beleg twee blini's met een beetje rucola en
        leg daar een plakje rosbief overheen. Schep er een theelepel pesto
        op en werk af met de Parmezaanse kaas.
  \item Voor de geitenkaasvariant: leg op twee blini's een schijfje geitenkaas
        en daarop een stukje zongedroogde tomaat, in reepjes gesneden of
        in zijn geheel. Werk af met een toefje honing en het takje tijm.
  \item Voor de appel-dillevariant: was en snijd de appel in piepkleine
        blokjes. Meng de zure room met de fijngesneden bieslook en een
        beetje peper. Schep op twee blini's een toef bieslookroom en leg
        daar een plakje gerookte zalm overheen. Werk af met de appelblokjes
        en dille. Serveer direct.
\end{steps}
\end{recipe}
